\documentclass[11pt, a4paper]{article}

% ── Codificação e idioma ──
\usepackage[utf8]{inputenc}
\usepackage[T1]{fontenc}
\usepackage[brazil]{babel}

% ── Layout ──
\usepackage[margin=1.8cm, top=2cm, bottom=1.8cm]{geometry}
\usepackage{enumitem}
\usepackage{multicol}

% ── Espaçamento ──
\setlength{\parskip}{2pt plus 1pt}
\setlength{\parindent}{0pt}

% ── Matemática ──
\usepackage{amsmath, amssymb}

% ── Tabelas e gráficos ──
\usepackage{booktabs}
\usepackage{array}
\usepackage{xcolor}
\usepackage{colortbl}
\usepackage[breakable]{tcolorbox}

% ── Cabeçalho e rodapé ──
\usepackage{fancyhdr}
\pagestyle{fancy}
\fancyhf{}
\lhead{\small Fundamentos de Sistemas Computacionais}
\rhead{\small Sistemas de Numera\c{c}\~{a}o}
\cfoot{\thepage}
\renewcommand{\headrulewidth}{0.4pt}

% ── Configurações de listas ──
\setlist[enumerate,1]{label=\textbf{\arabic*.}, leftmargin=*, itemsep=4pt, topsep=4pt, parsep=1pt}
\setlist[enumerate,2]{label=\alph*), leftmargin=1.5em, itemsep=1pt, topsep=2pt, parsep=0pt}
\setlist[itemize]{itemsep=1pt, topsep=2pt, parsep=0pt}

% ── Caixas de destaque ──
\newtcolorbox{dica}{
  colback=blue!5!white,
  colframe=blue!50!black,
  title={\textbf{Dica}},
  fonttitle=\small,
  sharp corners,
  boxrule=0.5pt,
  left=4pt, right=4pt, top=2pt, bottom=2pt,
  before skip=4pt, after skip=4pt
}

\newtcolorbox{contexto}{
  colback=green!5!white,
  colframe=green!50!black,
  sharp corners,
  boxrule=0.5pt,
  left=4pt, right=4pt, top=3pt, bottom=3pt,
  before skip=4pt, after skip=4pt
}

\newtcolorbox{curiosidade}{
  colback=yellow!5!white,
  colframe=yellow!60!black,
  title={\textbf{Curiosidade}},
  fonttitle=\small,
  sharp corners,
  boxrule=0.5pt,
  left=4pt, right=4pt, top=2pt, bottom=2pt,
  before skip=4pt, after skip=4pt
}

% ══════════════════════════════════════════
\begin{document}

\begin{center}
  {\LARGE \textbf{Lista de Exerc\'{i}cios}}\\[4pt]
  {\Large Sistemas de Numera\c{c}\~{a}o}\\[6pt]
  {\large 98800-04 -- Fundamentos de Sistemas Computacionais}\\[3pt]
  {\normalsize Prof.\ Ia\c{c}an\~{a} Ianiski Weber}
\end{center}

\vspace{2pt}
\hrule
\vspace{8pt}

\begin{enumerate}

\item Complete a frase: ``Com $N$~bits, \'{e} poss\'{i}vel representar no m\'{a}ximo \rule{2cm}{0.4pt}~coisas distintas.''
\item Escreva a \textbf{forma polinomial} de cada n\'{u}mero abaixo e calcule seu valor em decimal:
\begin{multicols}{3}
\begin{enumerate}
  \item $10110_2$
  \item $1101\,0011_2$
  \item $537_8$
  \item $1A3_{16}$
  \item $7F_{16}$
  \item $FF_{16}$
\end{enumerate}
\end{multicols}

\item Converta os seguintes n\'{u}meros \textbf{decimais para bin\'{a}rio}, usando divis\~{o}es sucessivas. Mostre todo o processo.
\begin{multicols}{3}
\begin{enumerate}
  \item $25_{10}$
  \item $100_{10}$
  \item $255_{10}$
  \item $197_{10}$
  \item $1024_{10}$
  \item $2023_{10}$
\end{enumerate}
\end{multicols}

\item Converta os seguintes n\'{u}meros \textbf{decimais para hexadecimal}. Voc\^{e} pode usar o m\'{e}todo que preferir (direto ou passando por bin\'{a}rio).
\begin{multicols}{3}
\begin{enumerate}
  \item $31_{10}$
  \item $200_{10}$
  \item $1000_{10}$
  \item $4095_{10}$
  \item $65535_{10}$
  \item $48879_{10}$
\end{enumerate}
\end{multicols}
\item Converta os seguintes n\'{u}meros \textbf{bin\'{a}rios para hexadecimal}, agrupando em \emph{nibbles}:
\begin{multicols}{2}
\begin{enumerate}
  \item $1011\,0110_2$
  \item $1111\,1111_2$
  \item $10\,1010\,1010_2$
  \item $1\,0000\,0000\,0000_2$
\end{enumerate}
\end{multicols}

\item Converta os seguintes n\'{u}meros \textbf{hexadecimais para bin\'{a}rio}:
\begin{multicols}{3}
\begin{enumerate}
  \item $\texttt{0xCA}$
  \item $\texttt{0xFE}$
  \item $\texttt{0xDEAD}$
  \item $\texttt{0xBEEF}$
  \item $\texttt{0xC0DE}$
  \item $\texttt{0xCAFE}$
\end{enumerate}
\end{multicols}

\begin{curiosidade}
Percebeu algo curioso nos n\'{u}meros hexadecimais acima? Programadores adoram criar ``palavras'' usando d\'{i}gitos hexadecimais (A--F). Esses n\'{u}meros s\~{a}o chamados de \emph{hexspeak}. Alguns exemplos famosos: \texttt{0xDEADBEEF}, \texttt{0xCAFEBABE} (primeiros bytes de todo arquivo Java \texttt{.class} -- usados pelo sistema para identificar o tipo do arquivo), \texttt{0xFEEDFACE} (idem para bin\'{a}rios Mach-O no macOS).
\end{curiosidade}

\item Responda:
\begin{multicols}{2}
\begin{enumerate}
  \item Quantos valores distintos um \emph{nibble} (4~bits) pode representar?
  \item Quantos valores distintos um \emph{byte} (8~bits) pode representar?
  \item Quantos d\'{i}gitos hexadecimais s\~{a}o necess\'{a}rios para representar um \emph{byte}?
  \item E para representar uma \emph{word} de 32~bits?
\end{enumerate}
\end{multicols}

\item Sem fazer c\'{a}lculos, determine rapidamente o valor decimal de:
\begin{multicols}{4}
\begin{enumerate}
  \item $2^{10}$
  \item $2^{16}$
  \item $2^{20}$
  \item $2^{32}$
\end{enumerate}
\end{multicols}
\vspace{-4pt}
\begin{dica}
Memorizar as pot\^{e}ncias de~2 \'{e} essencial em computa\c{c}\~{a}o. Lembre: $2^{10} = 1024 \approx 10^3$ (kilo), $2^{20} \approx 10^6$ (mega), $2^{30} \approx 10^9$ (giga).
\end{dica}
\item Qual \'{e} o n\'{u}mero m\'{i}nimo de bits necess\'{a}rios para representar os 12~meses do ano? Quantos valores adicionais poderiam ser representados sem bits extras?

\item A tabela ASCII usa 7~bits (arredondados para 1~byte na pr\'{a}tica). Quantos caracteres distintos 7~bits podem representar? Por que isso foi suficiente para o ingl\^{e}s, mas n\~{a}o para o portugu\^{e}s?

\item C\'{o}digos de cor HTML usam 24~bits (3~bytes), na forma \texttt{\#RRGGBB}, onde cada componente (R, G, B) usa 1~byte.
\begin{multicols}{2}
\begin{enumerate}
  \item Quantas cores distintas essa codifica\c{c}\~{a}o permite?
  \item O branco puro \'{e} \texttt{\#FFFFFF}. Qual \'{e} esse valor em decimal?
  \item O vermelho puro \'{e} \texttt{\#FF0000}. Qual o valor decimal de cada componente (R, G, B)?
  \item Invente uma cor, escreva seu c\'{o}digo hexadecimal e calcule o valor decimal equivalente.
\end{enumerate}
\end{multicols}

\item Um endere\c{c}o IPv4 usa 32~bits. Um endere\c{c}o IPv6 usa 128~bits.
\begin{enumerate}
  \item Quantos endere\c{c}os IPv4 existem? Expresse como pot\^{e}ncia de~2 e tamb\'{e}m em nota\c{c}\~{a}o decimal aproximada.
  \item Quantos endere\c{c}os IPv6 existem? Expresse como pot\^{e}ncia de~2.
  \item A popula\c{c}\~{a}o mundial \'{e} de aproximadamente $8 \times 10^9$ pessoas. Quantos endere\c{c}os IPv6 cada pessoa poderia ter, em m\'{e}dia? (Apresente o resultado como pot\^{e}ncia de~2 ou de~10.)
\end{enumerate}
\item Realize as seguintes \textbf{adi\c{c}\~{o}es em bin\'{a}rio} (unsigned, 8~bits). Indique se houve overflow.
\begin{multicols}{2}
\begin{enumerate}
  \item $0101\,0011_2 + 0010\,1100_2$
  \item $1100\,0011_2 + 0011\,1101_2$
  \item $1111\,0000_2 + 0001\,0001_2$
  \item $1011\,1011_2 + 0111\,0110_2$
\end{enumerate}
\end{multicols}

\item Realize as seguintes \textbf{subtra\c{c}\~{o}es em bin\'{a}rio} (8~bits):
\begin{multicols}{2}
\begin{enumerate}
  \item $1001\,0110_2 - 0011\,0101_2$
  \item $1111\,1111_2 - 0000\,0001_2$
  \item $1000\,0000_2 - 0000\,0001_2$
  \item $0101\,0101_2 - 0010\,1010_2$
\end{enumerate}
\end{multicols}
\vspace{-4pt}
\begin{dica}
$A - B = A + (-B)$
\end{dica}

\item Qual \'{e} o maior n\'{u}mero inteiro sem sinal (\emph{unsigned}) que pode ser representado com:
\begin{multicols}{3}
\begin{enumerate}
  \item 4~bits?
  \item 8~bits?
  \item 16~bits?
  \item 32~bits?
  \item $n$~bits?
\end{enumerate}
\end{multicols}

\item Um contador bin\'{a}rio de 4~bits come\c{c}a em $0000$ e incrementa de~1 em~1. Escreva a sequ\^{e}ncia completa de valores at\'{e} voltar a $0000$ (ap\'{o}s overflow). Quantos incrementos foram necess\'{a}rios?

\item Usando \textbf{complemento de dois com 8~bits}, represente os seguintes valores. Mostre o processo (complementar + somar~1):
\begin{multicols}{3}
\begin{enumerate}
  \item $+23$
  \item $-23$
  \item $+1$
  \item $-1$
  \item $+127$
  \item $-128$
\end{enumerate}
\end{multicols}

\item Converta os seguintes n\'{u}meros em \textbf{complemento de dois para decimal}:
\begin{multicols}{2}
\begin{enumerate}
  \item $0110_2$ \quad (4~bits)
  \item $1101_2$ \quad (4~bits)
  \item $0110\,1111_2$ \quad (8~bits)
  \item $1101\,1011\,0001\,1100_2$ \quad (16~bits)
\end{enumerate}
\end{multicols}
\vspace{-4pt}
\begin{dica}
Lembre-se que para o formato polinomial de um n\'{u}mero em complemento de dois o bit mais significativo tem peso negativo: $-2^{n-1} + 2^{n-2} + \ldots + 2^0$
\end{dica}

\newpage
\item Em complemento de dois com 12~bits:
\begin{multicols}{2}
\begin{enumerate}
  \item Qual \'{e} o maior n\'{u}mero positivo represent\'{a}vel? (em bin\'{a}rio e decimal)
  \item Qual \'{e} o n\'{u}mero negativo de maior magnitude? (em bin\'{a}rio e decimal)
\end{enumerate}
\end{multicols}
\item Generalize: em complemento de dois com $n$~bits:
\begin{enumerate}
  \item Qual \'{e} o maior positivo? Expresse como pot\^{e}ncia de~2.
  \item Qual \'{e} o menor negativo? Expresse como pot\^{e}ncia de~2.
  \item Quantos n\'{u}meros positivos, negativos e zeros existem?
\end{enumerate}

\item Expresse o valor $-22$ como inteiro em complemento de dois usando:
\begin{multicols}{3}
\begin{enumerate}
  \item 8~bits
  \item 16~bits
  \item 32~bits
\end{enumerate}
\end{multicols}
\vspace{-2pt}
Voc\^{e} percebe um padr\~{a}o? Esse fen\^{o}meno \'{e} chamado \textbf{extens\~{a}o de sinal} (\emph{sign extension}). Descreva-o em uma frase.

\item Explique por que, em complemento de dois com $n$~bits, o valor $-2^{n-1}$ \textbf{n\~{a}o} tem um inverso (positivo) represent\'{a}vel.
\item Descreva as condi\c{c}\~{o}es que indicam que ocorreu \textbf{overflow} ao somar dois n\'{u}meros em complemento de dois.

\item Os n\'{u}meros a seguir est\~{a}o em complemento de dois com \textbf{4~bits}. Realize as opera\c{c}\~{o}es, converta operandos e resultados para decimal e indique se houve overflow:
\begin{multicols}{4}
\begin{enumerate}
  \item $0011 + 1100$
  \item $0111 + 1111$
  \item $1110 + 1000$
  \item $0110 + 0010$
  \item $0101 + 0100$
  \item $1001 + 1010$
  \item $0111 + 0001$
  \item $1000 + 1111$
\end{enumerate}
\end{multicols}

\item Considere a soma em 4~bits (signed): $1100_2 + 0110_2$
\begin{enumerate}
  \item Interprete como \textbf{unsigned} e verifique se h\'{a} overflow.
  \item Interprete como \textbf{signed (complemento de dois)} e verifique se h\'{a} overflow.
  \item Explique por que a mesma opera\c{c}\~{a}o pode gerar overflow em uma interpreta\c{c}\~{a}o e n\~{a}o na outra.
\end{enumerate}

\item Um programador armazena a idade de uma pessoa em um inteiro com sinal de 8~bits (complemento de dois). Qual \'{e} a maior idade que pode ser armazenada? O que aconteceria se ele tentasse armazenar a idade de Jeanne Calment, que viveu 122~anos?

\item \textbf{O Enigma do Od\^{o}metro Bin\'{a}rio:} Imagine um od\^{o}metro de carro que, em vez de d\'{i}gitos decimais, usa 16~bits bin\'{a}rios (unsigned).
\begin{enumerate}
  \item Qual a quilometragem m\'{a}xima que esse od\^{o}metro pode registrar?
  \item Se o carro roda 50~km por dia, ap\'{o}s quantos dias o od\^{o}metro ``zera'' (overflow)?
  \item Se o od\^{o}metro usasse 32~bits, quantos anos levariam para zerar, nas mesmas condi\c{c}\~{o}es?
\end{enumerate}

\item \textbf{Mensagem Secreta:} Uma mensagem foi codificada convertendo cada letra para sua posi\c{c}\~{a}o no alfabeto (A=1, B=2, \ldots, Z=26) e depois para bin\'{a}rio com 5~bits. Decodifique: \texttt{01000 01001 01110 00001 10010 01001 01111}

\item \textbf{O Bug do Pac-Man:} No jogo original \emph{Pac-Man} (1980), o contador de n\'{i}vel era armazenado em um \'{u}nico byte (8~bits, unsigned). Quando o jogador chegava ao n\'{i}vel~256, ocorria um famoso \emph{bug}.
\begin{enumerate}
  \item Explique, do ponto de vista de sistemas de numera\c{c}\~{a}o, o que aconteceu no n\'{i}vel~256.
  \item Qual seria o maior n\'{i}vel poss\'{i}vel se o contador usasse 16~bits?
  \item E se usasse um inteiro com sinal de 8~bits (complemento de dois)?
\end{enumerate}

\newpage
\item \textbf{Nuclear Gandhi:} Na s\'{e}rie de jogos \emph{Civilization}, existe uma lenda urbana de que o l\'{i}der Gandhi, normalmente pac\'{i}fico, se tornava extremamente agressivo. A ``agressividade'' de Gandhi come\c{c}ava em~1 (de~0 a~255, armazenada em 1~byte unsigned). Ao adotar a democracia, a agressividade era reduzida em~2.
\begin{enumerate}
  \item Calcule $1 - 2$ em aritm\'{e}tica de 8~bits unsigned. Qual valor resultaria?
  \item Explique o comportamento inesperado usando o conceito de overflow/underflow.
  \item Se a vari\'{a}vel fosse um inteiro com sinal de 8~bits, qual seria o resultado de $1-2$? O bug aconteceria?
\end{enumerate}
\item \textbf{O Bug do Ano 2038:} Muitos sistemas Unix armazenam timestamps como o n\'{u}mero de segundos desde 1\textordmasculine{} de janeiro de 1970 (``\emph{Unix epoch}''), usando um inteiro com sinal de 32~bits (complemento de dois).
\begin{enumerate}
  \item Qual \'{e} o maior valor positivo que um inteiro com sinal de 32~bits pode armazenar?
  \item Quantos segundos h\'{a} em um ano (considere 365,25 dias)?
  \item Em que ano aproximado o contador ir\'{a} ``estourar'' (overflow)?
  \item Qual \'{e} a solu\c{c}\~{a}o adotada pelos sistemas modernos?
\end{enumerate}

\item \textbf{Endere\c{c}os de Mem\'{o}ria:} Um processador usa endere\c{c}os de 32~bits para acessar a mem\'{o}ria RAM.
\begin{dica}
A mem\'{o}ria \'{e} endere\c{c}ada a byte: cada endere\c{c}o aponta para exatamente 1~byte.
\end{dica}
\vspace{-2pt}
\begin{enumerate}
  \item Quantos bytes de mem\'{o}ria podem ser endere\c{c}ados?
  \item Converta essa quantidade para GiB ($2^{30}$~bytes).
  \item Explique por que sistemas operacionais de 32~bits t\^{e}m um limite pr\'{a}tico de~4~GiB de RAM.
  \item Processadores de 64~bits resolvem esse problema. Quantos bytes poderiam, em teoria, ser endere\c{c}ados? \'{E} um valor pr\'{a}tico?
\end{enumerate}

\item \textbf{Placas de Ve\'{i}culos:} As placas de ve\'{i}culos no Brasil (padr\~{a}o Mercosul) t\^{e}m o formato \texttt{LLLNLNN}, onde L~=~letra (26~possibilidades) e N~=~d\'{i}gito (10~possibilidades).
\begin{enumerate}
  \item Quantas placas distintas esse formato permite?
  \item Quantos bits seriam necess\'{a}rios para codificar cada placa como um n\'{u}mero bin\'{a}rio \'{u}nico?
\end{enumerate}

\item \textbf{Pixel Art:} Uma imagem em escala de cinza usa 8~bits por pixel (0~=~preto, 255~=~branco).
\begin{enumerate}
  \item Quantos tons de cinza distintos essa codifica\c{c}\~{a}o permite?
  \item Considere uma imagem de $1920 \times 1080$ pixels. Quantos bytes s\~{a}o necess\'{a}rios para armazen\'{a}-la (sem compress\~{a}o)?
  \item Se cada pixel usasse 24~bits (cor RGB), quantos bytes seriam necess\'{a}rios?
  \item Converta as respostas anteriores para MiB ($2^{20}$~bytes).
\end{enumerate}
\item \textbf{Multiplica\c{c}\~{a}o e divis\~{a}o por pot\^{e}ncias de~2:} Em bin\'{a}rio, multiplicar por~2 equivale a deslocar todos os bits uma posi\c{c}\~{a}o \`{a} esquerda (\emph{shift left}), e dividir por~2 equivale a deslocar uma posi\c{c}\~{a}o \`{a} direita (\emph{shift right}).
\begin{enumerate}
  \item Dado $0010\,1100_2$ (44 em decimal), calcule o resultado de desloc\'{a}-lo 1~bit \`{a} esquerda. Verifique convertendo para decimal.
  \item Dado $0110\,0000_2$ (96 em decimal), calcule o resultado de desloc\'{a}-lo 2~bits \`{a} direita. Verifique.
  \item Generalize: deslocar $k$~bits \`{a} esquerda equivale a multiplicar por quanto? E deslocar $k$~bits \`{a} direita?
  \item Um programador precisa multiplicar um n\'{u}mero por~10. Usando apenas shifts e adi\c{c}\~{o}es, como ele pode fazer isso? (Dica: $10 = 8 + 2$.)
\end{enumerate}

\newpage
\item Converta os seguintes n\'{u}meros bin\'{a}rios fracion\'{a}rios para decimal:
\begin{contexto}
Assim como no sistema decimal temos casas ap\'{o}s a v\'{i}rgula ($10^{-1}, 10^{-2}, \ldots$), em bin\'{a}rio as posi\c{c}\~{o}es ap\'{o}s o ponto bin\'{a}rio valem $2^{-1} = 0.5$, $2^{-2} = 0.25$, $2^{-3} = 0.125$, etc.
\end{contexto}
\vspace{-2pt}
\begin{multicols}{3}
\begin{enumerate}
  \item $0.1_2$
  \item $0.11_2$
  \item $0.101_2$
  \item $1.1_2$
  \item $10.01_2$
  \item $11.111_2$
\end{enumerate}
\end{multicols}

\item Converta os seguintes n\'{u}meros decimais fracion\'{a}rios para bin\'{a}rio (use at\'{e} 6~casas bin\'{a}rias):
\begin{multicols}{3}
\begin{enumerate}
  \item $0.5_{10}$
  \item $0.75_{10}$
  \item $0.625_{10}$
  \item $0.1_{10}$
  \item $0.3_{10}$
  \item $0.7_{10}$
\end{enumerate}
\end{multicols}

\item Voc\^{e} notou que $0.1_{10}$ n\~{a}o tem uma representa\c{c}\~{a}o bin\'{a}ria exata (finita)?
\begin{enumerate}
  \item Execute o seguinte teste mental (ou em um computador): em Python, $0.1 + 0.2 == 0.3$? O resultado te surpreende?
  \item Que implica\c{c}\~{o}es pr\'{a}ticas isso tem para softwares financeiros?
\end{enumerate}

\item Calcule os resultados das seguintes opera\c{c}\~{o}es bit a bit (8~bits):
\begin{contexto}
Al\'{e}m de opera\c{c}\~{o}es aritm\'{e}ticas, computadores realizam opera\c{c}\~{o}es \textbf{l\'{o}gicas} bit a bit: AND (\&), OR ($|$), XOR (\texttt{\^{}}) e NOT ($\sim$). Cada opera\c{c}\~{a}o \'{e} aplicada individualmente a cada par de bits correspondentes.
\end{contexto}
\vspace{-2pt}
\begin{multicols}{2}
\begin{enumerate}
  \item $1010\,0110$ AND $1100\,0011$
  \item $1010\,0110$ OR $1100\,0011$
  \item $1010\,0110$ XOR $1100\,0011$
  \item NOT $1010\,0110$
\end{enumerate}
\end{multicols}

\item \textbf{M\'{a}scaras de bits (\emph{bitmasks}):} M\'{a}scaras de bits s\~{a}o usadas para extrair, definir ou limpar bits espec\'{i}ficos de um n\'{u}mero.
\begin{enumerate}
  \item Dado o byte $\texttt{0xB7}$, aplique a m\'{a}scara $\texttt{0x0F}$ com AND. O que essa opera\c{c}\~{a}o extrai?
  \item Dado o byte $\texttt{0xB7}$, aplique a m\'{a}scara $\texttt{0xF0}$ com AND e desloque 4~bits \`{a} direita. O que voc\^{e} obt\'{e}m?
  \item Explique como usar OR para ``ligar'' o bit~3 (contando de~0, da direita) de um byte, sem alterar os demais bits.
  \item Explique como usar AND para ``desligar'' o bit~5 de um byte, sem alterar os demais.
  \item Uma propriedade do XOR: $A$ \texttt{\^{}} $B$ \texttt{\^{}} $B$ $= A$. Explique por que essa propriedade \'{e} \'{u}til em criptografia simples e em algoritmos de troca de vari\'{a}veis sem usar uma vari\'{a}vel tempor\'{a}ria.
\end{enumerate}
\item \textbf{Swap sem vari\'{a}vel tempor\'{a}ria}

Dados $A = 0101\,1010$ e $B = 1100\,0011$, execute a sequ\^{e}ncia:
\begin{itemize}[label=$\bullet$]
  \item \texttt{A = A \^{} B}
  \item \texttt{B = A \^{} B}
  \item \texttt{A = A \^{} B}
\end{itemize}
\vspace{-2pt}
Mostre o valor de $A$ e $B$ ap\'{o}s cada passo. O que aconteceu?

\item \textbf{O Gangnam Style e o YouTube:} Em 2014, o clipe de ``Gangnam Style'' do PSY ultrapassou 2.147.483.647 visualiza\c{c}\~{o}es no YouTube. O contador do YouTube parou de funcionar.
\begin{enumerate}
  \item Reconhece esse n\'{u}mero? Qual \'{e} sua rela\c{c}\~{a}o com pot\^{e}ncias de~2?
  \item O YouTube usava um inteiro com sinal de 32~bits para contar visualiza\c{c}\~{o}es. Qual foi o problema?
  \item Ap\'{o}s o incidente, o YouTube migrou para um inteiro com sinal de 64~bits. Qual \'{e} o novo limite? \'{E} suficiente?
\end{enumerate}

\newpage
\item \textbf{C\'{o}digos de Barras e QR Codes:} Um c\'{o}digo de barras EAN-13 (usado em produtos) codifica 13~d\'{i}gitos decimais.
\begin{enumerate}
  \item Quantas combina\c{c}\~{o}es distintas de 13~d\'{i}gitos decimais existem?
  \item Quantos bits seriam necess\'{a}rios para representar esse mesmo n\'{u}mero de combina\c{c}\~{o}es em bin\'{a}rio puro?
  \item Compare com a codifica\c{c}\~{a}o BCD: quantos bits seriam necess\'{a}rios em BCD para 13~d\'{i}gitos?
\end{enumerate}

\item \textbf{Capacidade de Armazenamento ao Longo da Hist\'{o}ria:} O primeiro disco r\'{i}gido da IBM (1956, modelo 305 RAMAC) armazenava 5~MB. Hoje, discos de 20~TB s\~{a}o comuns.
\begin{enumerate}
  \item Quantos bits s\~{a}o necess\'{a}rios para representar 5~MB e 20~TB?
  \item Quantas vezes a capacidade de armazenamento aumentou de 1956 at\'{e} hoje? Expresse como pot\^{e}ncia de~2 aproximada.
\end{enumerate}

\end{enumerate}

\vspace{8pt}
\hrule
\vspace{8pt}

% ══════════════════════════════════════════
\section*{Desafios}
% ══════════════════════════════════════════

\begin{enumerate}[resume]

\item \textbf{O Curioso Caso do $-0$ no Espa\c{c}o}

Em 1996, o foguete Ariane~5 da Ag\^{e}ncia Espacial Europeia explodiu 37~segundos ap\'{o}s o lan\c{c}amento. A causa raiz foi a convers\~{a}o de um n\'{u}mero de ponto flutuante de 64~bits para um inteiro com sinal de 16~bits, gerando overflow.
\begin{enumerate}
  \item Qual \'{e} o intervalo de valores de um inteiro com sinal de 16~bits (complemento de dois)?
  \item Se a velocidade horizontal do foguete era armazenada como um valor que excedia $32\,767$, o que aconteceu na convers\~{a}o?
  \item Que li\c{c}\~{a}o sobre sistemas de numera\c{c}\~{a}o e overflow esse acidente ensina?
\end{enumerate}

\item \textbf{Paradoxo do Anivers\'{a}rio em Hashes:} Fun\c{c}\~{o}es hash mapeiam dados de tamanho arbitr\'{a}rio para um valor de tamanho fixo. O ``paradoxo do anivers\'{a}rio'' implica que colis\~{o}es s\~{a}o prov\'{a}veis ap\'{o}s aproximadamente $\sqrt{2^n}$ tentativas, onde $n$~\'{e} o n\'{u}mero de bits do hash.
\begin{enumerate}
  \item Para um hash de 32~bits, ap\'{o}s quantas tentativas uma colis\~{a}o se torna prov\'{a}vel? Expresse como pot\^{e}ncia de~2.
  \item Para um hash de 256~bits (SHA-256), calcule o mesmo valor.
  \item A rede Bitcoin processa $\sim 10^{20}$ hashes por ano. Quantos anos seriam necess\'{a}rios para encontrar uma colis\~{a}o em SHA-256?
\end{enumerate}

\end{enumerate}

\vfill
\begin{center}
\textit{``Existem 10 tipos de pessoas no mundo: as que entendem bin\'{a}rio e as que n\~{a}o entendem.''}
\end{center}

\end{document}
